\session{6}{10/24}{1限}{2}{組織・人材とAI協働}

\goals{
  \item 文書作成・会議・企画といった知識労働が生成AIでどう変わるかを具体的に説明できる。
  \item 人とAIの協働の形と、求められるスキルの変化を理解する。
  \item AI導入に対する組織の抵抗の要因と、変革管理の基本を説明できる。
}

\guidesec{解説}

\topic{知識労働の変化}
生成AIが最も直接的に影響しているのは、文書を書く、会議をする、企画を立てるといった知識労働です。
\par\noindent\begin{gtabh}{colspec={Q[l,wd=20mm] X[l] X[l]}}
  仕事 & AIが担えるようになったこと & 人に残ること \\
  文書作成 & 下書き、要約、翻訳、体裁の統一、誤字の指摘 & 何を伝えるかの決定、事実の確認、相手に応じた判断、最終責任 \\
  会議 & 議事録の作成、論点の整理、決定事項の抽出、資料の要約 & 議論そのもの、合意の形成、決定 \\
  企画 & アイデアの大量生成、競合や市場の情報整理、企画書の構成案 & 目的の設定、案の選択、実行の意思決定、関係者の説得 \\
  分析 & データの集計、グラフ作成、仮説の候補出し & 問いの設定、結果の解釈、意思決定への反映（第8〜11回） \\
\end{gtabh}
共通するのは、「作る」作業の多くをAIが担い、人の仕事が「決める」「確かめる」「責任を持つ」に寄っていくことです。
生産性の研究では、生成AIは経験の浅い担当者の成果を大きく引き上げ、熟練者との差を縮める傾向が報告されています。
一方で、AIの誤りに気づけるかどうかは専門知識に依存するため、専門性の価値がなくなるわけではありません。

\topic{人とAIの協働}
人とAIの協働は、AIの提案を人が確認・修正しながら成果を出す働き方です。AIをどの役割で使うかを意識すると、協働を設計しやすくなります。
\begin{itemize}
  \item \textbf{下書き担当}　AIが最初の案を作り、人が修正して仕上げる。速さを得る使い方。
  \item \textbf{壁打ち相手}　人の案に対してAIが質問・反論・別案を出す。質を上げる使い方。
  \item \textbf{点検役}　人が書いたものの誤り、抜け、論理の飛躍をAIに指摘させる。
\end{itemize}
「AIが作って人が確認する」だけでなく、「人が作ってAIが確認する」方向にも使えることを覚えておいてください。
また、どの使い方でも、AIの寄与と自分の判断を区別して示す（AI Policy）ことが協働の前提です。

\topic{求められるスキルの変化}
AIが文章や案を出せる時代には、次の力の重要性が高まります。
\begin{itemize}
  \item \textbf{問いを立てる力}　何を知りたいか、何を解決したいかを言葉にする力。プロンプトの質はここで決まる。
  \item \textbf{評価する力}　出力が正しいか、目的に合っているか、使ってよいかを判断する力。対象領域の知識が必要。
  \item \textbf{説明する力}　AIの出力を根拠とともに他者に説明し、責任を持って決める力。
  \item \textbf{学び続ける力}　道具の変化が速いため、新しい使い方を試し、更新し続ける力。
\end{itemize}
これらは、従来「経験を積んだ管理者」に求められていた力です。AIによって、若手にも早くからこの力が求められるようになったと言えます。

\topic{組織の抵抗と変革管理}
AI導入に対して、組織では抵抗が生じます。主な要因は、仕事がなくなる不安、使い方への不慣れ、今までのやり方への愛着、
評価や責任の不明確さです。抵抗は非合理なものではなく、合理的な不安の表れとして理解する必要があります。

変革を定着させる管理（チェンジマネジメント）の基本は次の通りです。
\begin{itemize}
  \item \textbf{目的の共有}　なぜ導入するのか、従業員にとって何が良くなるのかを、経営層が自分の言葉で伝える。
  \item \textbf{段階的な導入}　小さな範囲で試し、成果と課題を示してから広げる。短期的な成果を見えるようにする。
  \item \textbf{教育と支援}　使い方の研修、相談窓口、使ってよい範囲のルール（第13回）を用意する。
  \item \textbf{現場の参加}　導入の設計に現場の担当者を巻き込む。使う人が決めた道具は使われる。
  \item \textbf{評価と責任の明確化}　AIを使った成果をどう評価するか、誤りが起きたときの責任をどう扱うかを決める。
\end{itemize}
コッターの変革8段階（危機意識の醸成、推進チームの結成、ビジョンの策定と共有、短期的成果の実現、定着）も参考になります。
経営層の役割は、技術の詳細を理解することではなく、目的と方針を示し、体制と投資を決め、責任を持って推進することです。

\topic{キーワード}
\par\noindent\begin{keywords}{}
  知識労働 & 情報を扱い、判断や創造を伴う仕事。文書作成、会議、企画、分析など。 \\
  人とAIの協働 & AIの提案を人が確認・修正しながら成果を出す働き方。下書き・壁打ち・点検の3つの役割。 \\
  問いを立てる力・評価する力 & AI時代に重要性が高まるスキル。 \\
  組織の抵抗 & 雇用不安、不慣れ、愛着、責任の不明確さから生じる、変化への反発。 \\
  変革管理 & 目的の共有、段階的導入、教育・支援、現場の参加で変革を定着させる管理。 \\
\end{keywords}

\guidesec{演習課題}

\exercise{1}{AI演習：AI導入に賛成・反対の立場をAIに演じさせ、論点を整理する}
\instr{「ある中堅企業（従業員300人、事務部門80人）が、事務部門に生成AIを全面導入する」という場面を想定し、次のプロンプトで賛成派と反対派の議論をAIに行わせてください。出力から論点を抜き出して表にし、各論点について「自分ならどう判断するか」と、その根拠を書いてください。}
\begin{promptbox}[賛成・反対の討論]
従業員300人の中堅企業が、事務部門（80人）に生成AIを全面導入しようとしています。あなたは2人の役員を演じてください。Aは導入推進派の経営企画担当役員、Bは慎重派の管理部門担当役員です。AとBが交互に発言する形で、5往復の議論を書いてください。各発言は具体的な業務の例と懸念を含めてください。最後に、議論に出た論点を箇条書きで整理してください。
\end{promptbox}

\exercise{2}{自分のスキルの棚卸し}
\instr{「問いを立てる力」「評価する力」「説明する力」「学び続ける力」について、自分の現状を5段階で自己評価し、それぞれについて今後半年で取り組むことを1つ書いてください。}

\exercise{3}{変革計画を立てる}
\instr{演習1の企業の経営者として、事務部門へのAI導入を1年で定着させる計画を、「目的の共有」「段階的な導入」「教育と支援」「現場の参加」「評価と責任」の5項目について、それぞれ具体的な施策として書いてください。}

\guidesec{模範解答}

\begin{answer}{演習1}
\par\noindent\begin{gtabh}{colspec={Q[l,wd=26mm] X[l] X[l] X[l,wd=40mm]}}
  論点 & 推進派（A） & 慎重派（B） & 自分の判断 \\
  生産性 & 文書作成と会議の時間を3割減らせる & 確認の手間が増え、差し引きの効果は不明 & 対象業務を絞って効果を測ってから判断する（第7回のKPI）。 \\
  雇用 & 削減ではなく、付加価値の高い仕事への配置転換 & 従業員は削減と受け取り、士気が下がる & 配置転換の具体像を先に示す。削減しない方針なら明言する。 \\
  情報漏えい & 法人契約のツールなら入力は学習に使われない & 個人のアカウントで使う人が出る（シャドーAI） & 公式ツールの提供とルール整備を導入と同時に行う。 \\
  品質・責任 & 人が確認すれば問題ない & 誰が確認したかが曖昧になる & 成果物に確認者を明記し、AI活用記録の形式を決める。 \\
  教育 & 使えば慣れる & 年齢や部署で習熟に差が出る & 部署ごとに推進担当を置き、研修と相談窓口を設ける。 \\
\end{gtabh}
気づき：AIに両方の立場を演じさせると、自分が思いつかなかった反対意見（確認の手間、習熟の差）が出てくる。
ただし、AIの議論は一般論になりがちで、自社の具体的な状況（どの業務に何時間かかっているか）は自分で補う必要がある。
\end{answer}

\begin{answer}{演習2}
例：
\begin{itemize}
  \item 問いを立てる力：3。レポートのテーマは決められるが、問いの形（なぜ〜か）にするのに時間がかかる。取り組み：AI利用ガイドの復習教材「問いの立て方」に取り組む。
  \item 評価する力：2。AIの出力の誤りに気づけないことが多い。取り組み：出力の固有名詞と数値を毎回一次資料で確かめる習慣をつける。
  \item 説明する力：3。根拠を示して説明するのが苦手。取り組み：演習の答えに必ず「根拠」の欄を作る。
  \item 学び続ける力：4。新しいツールは試す方だが、記録していない。取り組み：試したツールと使い方を月に1回まとめる。
\end{itemize}
\end{answer}

\begin{answer}{演習3}
\begin{itemize}
  \item 目的の共有：社長が全社集会で「事務作業の時間を顧客対応と改善活動に振り向けるため」と目的を説明し、人員削減をしない方針を明言する。
  \item 段階的な導入：最初の3か月は経理と総務の各2名で、議事録と社内文書の下書きに限定して試行。効果（作業時間）と課題を全社に報告してから、6か月目に事務部門全体へ広げる。
  \item 教育と支援：全員に2時間の基礎研修、希望者に実践研修。部門ごとに推進担当を1名置き、チャットで質問できる窓口を作る。
  \item 現場の参加：試行メンバーが「使ってよい業務」と「使わない業務」のリストを作り、ルールの原案とする。
  \item 評価と責任：AIを使った成果物には確認者を明記する。誤りが起きた場合は個人を責めず、プロセスを見直す方針を示す。半年ごとに作業時間と満足度を測る。
\end{itemize}
\end{answer}

\guidesec{出席確認クイズの解説}
講義サイトの出席確認ページ（第6回）の5問です。
% 出席確認クイズ 第06回　組織・人材とAI協働
%   attendance/attendance06.html から生成（解説は本ガイド用に加筆）。

\quizq{1}{「人とAIの協働」の説明として適切なものはどれでしょうか?}%
  {\correct{AIの提案を人が確認・修正しながら成果を出す働き方}}%
  {AIがすべてを決め、人は関与しない}%
  {AI同士が競争する}%
  {人がAIをまったく使わない}%
  {正解は (a)。人とAIの協働では、AI の提案を人が確認・修正し、それぞれの得意分野を活かして成果を出します。(b) は全自動化であり、協働ではありません。AIを下書き担当・壁打ち相手・点検役のどれとして使うかを意識すると、協働の設計がしやすくなります。}

\quizq{2}{AI導入の際に組織で起こりやすい抵抗の要因として適切なものはどれでしょうか?}%
  {電力不足}%
  {\correct{仕事がなくなる不安や使い方への不慣れ}}%
  {為替の変動}%
  {天候の変化}%
  {正解は (b)。AI導入への抵抗は、雇用への不安と新しい道具への不慣れが主な要因です。(a)(c)(d) は外部環境の要因で、組織内部の抵抗とは関係がありません。抵抗は非合理ではなく、合理的な不安の表れとして理解する必要があります。}

\quizq{3}{変革管理(チェンジマネジメント)で重要なこととして適切なものはどれでしょうか?}%
  {罰則で強制する}%
  {一度にすべてを切り替える}%
  {\correct{目的の共有と段階的な導入、教育や支援}}%
  {導入を従業員に秘密にする}%
  {正解は (c)。変革を定着させるには、目的の共有、段階的な導入、教育・支援による不安の解消が重要です。(a)(b)(d) はいずれも抵抗を強めます。コッターの変革8段階（危機感の醸成、ビジョンの共有、短期的成果など）も参考になります。}

\quizq{4}{生成AIの普及によって重要性が高まると考えられるスキルはどれでしょうか?}%
  {電話応対のみ}%
  {\correct{適切な問いを立て、AIの出力を評価する力}}%
  {暗算の速さ}%
  {手書きの速さ}%
  {正解は (b)。AIが文章や案を出せる時代には、適切な問いを設計し、出力の質を評価・判断する力がより重要になります。(a)(c)(d) は機械に置き換えられやすい技能です。評価には対象領域の専門知識が必要で、専門性の価値がなくなるわけではありません。}

\quizq{5}{AI導入において経営層に求められる役割として適切なものはどれでしょうか?}%
  {現場に任せて一切関与しない}%
  {技術の詳細をすべて自分で実装する}%
  {\correct{目的と方針を示し、責任を持って推進する}}%
  {AIの利用をすべて禁止する}%
  {正解は (c)。経営層には、AI活用の目的と方針を示し、体制や投資を決め、責任を持って推進する役割があります。(b) の技術実装は専門家の仕事で、(a) の現場への丸投げは第7回で扱う失敗要因の1つです。(d) は競争力の低下とシャドーAIを招きます。}

